% Generated draft: requires mathematical and source review.
\documentclass[11pt]{article}
\usepackage[margin=25mm]{geometry}
\usepackage{fontspec}
\setmainfont{Latin Modern Roman}
\usepackage{amsmath,amssymb,mathtools,unicode-math}
\setmathfont{Latin Modern Math}
\usepackage{enumitem,needspace,xurl,hyperref}
\setlength{\parindent}{0pt}
\setlength{\parskip}{5pt}
\setlength{\emergencystretch}{3em}
\newcommand{\N}{\mathbb{N}}
\newcommand{\Z}{\mathbb{Z}}
\newcommand{\Q}{\mathbb{Q}}
\newcommand{\R}{\mathbb{R}}
\newcommand{\C}{\mathbb{C}}
\newcommand{\F}{\mathbb{F}}
\newcommand{\A}{\mathbb{A}}
\newcommand{\PP}{\mathbb{P}}
\newcommand{\E}{\mathbb{E}}
\newcommand{\eps}{\varepsilon}
\newcommand{\dd}{\,\mathrm{d}}
\begin{document}
\section*{UnsolvedMath ID: 30005619}
\textbf{Source problem code:} catalogue entry
\subsection{Definitions and mathematical statement}
Let \((M,g)\) be a fixed closed Riemannian surface, where closed means compact without boundary, and let \(i>1\) be an integer. The supplied record does not specify connectedness, the admissible regularity of magnetic potentials, the spectral indexing convention, or the normalization denoted by the bar.

A conventional interpretation, requiring confirmation against the source, takes a real magnetic potential to be a smooth real differential one-form \(A\in\Omega^1(M;\R)\), and defines the magnetic Laplacian on complex-valued functions by
\[
 d_Au=\dd u+\mathrm{i}uA,\qquad \Delta_{g,A}=d_A^*d_A,
\]
where the adjoint uses the metric \(g\) and its Riemannian volume measure. Under one-based indexing, its eigenvalues, repeated according to multiplicity, are written
\[
0\leq\lambda_1(M,g,A)\leq\lambda_2(M,g,A)\leq\cdots.
\]
A conventional surface normalization is
\[
\overline\lambda_j(M,g,A)=\operatorname{Area}_g(M)\lambda_j(M,g,A).
\]
These conventions are explanatory candidates; they are not definitions verified from the supplied excerpt.

The principal requested statement is that, for a fixed \((M,g)\) and a fixed integer \(i>1\),
\[
\inf_A\overline\lambda_i(M,g,A)\geq C(M,g)>0,
\]
with the infimum taken over all admissible real magnetic potentials. Under the provisional conventions above, the infimum would be over \(\Omega^1(M;\R)\). The constant must be independent of \(A\). Although the source writes \(C(M,g)\), it does not explicitly resolve whether the constant must also be independent of \(i\). Connectedness and the spectral conventions must be clarified before treating this as a fully specified conjecture.

The record additionally asks for a natural construction of critical metric-potential pairs \((g,A)\). It supplies neither the functional whose critical points are intended nor the allowed variations, constraints, or definition of criticality. Consequently this additional request cannot yet be formulated as a precise mathematical assertion.
\subsection{Short English statement}
For a fixed closed Riemannian surface and an eigenvalue index \(i>1\), seek a positive lower bound for the normalized magnetic eigenvalue that holds for every admissible real magnetic potential. Also seek a natural construction of critical metric-potential pairs. The normalization, indexing, connectedness assumptions, admissible potentials, and meaning of criticality require confirmation.
\subsection{Sources}
\begin{itemize}
\item Local UnsolvedMath registry: supplied complete source record, permanent ID \(30005619\), exact source problem code \texttt{catalogue entry}, title \emph{Magnetic-Laplacian Eigenvalue Infima on Closed Surfaces}. The record belongs to the Oberwolfach Reports Open Problems collection; no individual proposer is supplied.
\item Supplied source citation: \emph{Geometric Spectral Theory (2024)}, Oberwolfach Reports, DOI 10.4171/owr/2023/36, pp.\ 2017--2106. Supplied source URL and literature reference labeled ``original OWR report'': \url{https://doi.org/10.4171/owr/2023/36}. The record separately gives proposed year 2023 and the tag \texttt{owr-year:2023}; these are retained without reconciling them with the citation's 2024 label.
\item Literature assessment dated exactly \texttt{2026-08-22}: the supplied curation labels the problem open and states that available magnetic spectral results do not establish the fixed-surface lower bound in the stated generality. This is supplied dated curation, not an independently verified literature assessment. The background explicitly requests independent verification of time-sensitive status and cited claims.
\item The supplied statement status is \texttt{corrected\_verified}, with the supplied explanation that the statement was reconstructed and checked against the official Oberwolfach report. That verification is reported as source metadata; it has not been independently reproduced here.
\end{itemize}
\end{document}
